
\documentclass[a4paper,12pt]{article}

\input{glyphtounicode}
\pdfgentounicode=1

\usepackage[T1]{fontenc}
\usepackage[utf8]{inputenc}
\usepackage[english]{babel}
\usepackage[version=4]{mhchem}
\usepackage[dvipsnames]{xcolor}
\usepackage[backend=biber,style=ieee]{biblatex}
\usepackage{newtxtext,newtxmath}
\usepackage{microtype}
\usepackage{setspace}
\usepackage{titlesec}
\usepackage{tocloft}
\usepackage{indentfirst}
\usepackage{enumitem}
\usepackage{fancyhdr}
\usepackage{hyperref}
\usepackage{tabularx}
\usepackage{array}
\usepackage{makecell}
\usepackage{booktabs}
\usepackage{graphicx}
\usepackage{amsmath}
\usepackage{footmisc}
\usepackage{csquotes}
\usepackage{fontawesome5}
\usepackage[
  a4paper,
  top        = 2.00cm,
  bottom     = 2.00cm,
  left       = 2.00cm,
  right      = 2.00cm,
  headheight = 15pt,
  headsep    = 5pt,
  footskip   = 15pt
]{geometry}

\newcommand{\degreeawardvalue}{}
\newcommand{\universityvalue}{}
\newcommand{\addressvalue}{}
\newcommand{\unilogovalue}{}
\newcommand{\copyyearvalue}{}
\newcommand{\defenddatevalue}{}
\newcommand{\orcidvalue}{}
\newcommand{\rightsstatementvalue}{}
\newcommand{\degreeaward}[1]{\def\degreeawardvalue{#1}}
\newcommand{\university}[1]{\def\universityvalue{#1}}
\newcommand{\address}[1]{\def\addressvalue{#1}}
\newcommand{\unilogo}[1]{\def\unilogovalue{#1}}
\newcommand{\copyyear}[1]{\def\copyyearvalue{#1}}
\newcommand{\defenddate}[1]{\def\defenddatevalue{#1}}
\newcommand{\orcid}[1]{\def\orcidvalue{#1}}
\newcommand{\rightsstatement}[1]{\def\rightsstatementvalue{#1}}
\newcommand{\titleskip}{\vskip 4\bigskipamount}
\newcommand{\authorskip}{\vskip 2\bigskipamount}
\newcommand{\resumesection}[2]{\section{\textcolor{red}{#1}#2}}
\newcommand{\resumeSubHeadingListStart}{\begin{itemize}[leftmargin=0.0in, label={}, itemsep=2pt]}
\newcommand{\resumeSubHeadingListEnd}{\end{itemize}}
\newcommand{\resumeItemListStart}{\begin{itemize}}
\newcommand{\resumeItemListEnd}{\end{itemize}}
\newcommand{\resumeItem}[1]{
    \item\small{#1}
}
\newcommand{\resumeSubheading}[4]{
    \item
    \begin{tabular*}{1.0\textwidth}[t]{l@{\extracolsep{\fill}}r}
        \textbf{#1} & \textbf{\small #2} \\
        \small#3 & \small #4
    \end{tabular*}
}

\renewcommand{\cftsecleader}{\cftdotfill{\cftdotsep}}
\renewcommand{\cftsecpresnum}{Chapter~}
\renewcommand{\cftsecaftersnum}{:}
\renewcommand{\headrulewidth}{0pt}
\renewcommand{\footrulewidth}{0pt}
\renewcommand\labelitemi{$\vcenter{\hbox{\tiny$\bullet$}}$}
\renewcommand\labelitemii{$\vcenter{\hbox{\tiny$\bullet$}}$}

\pagestyle{fancy}
\fancyhf{}
% \fancyfoot[R]{\thepage}
\urlstyle{same}
\addbibresource{references.bib}
\hypersetup{
  colorlinks=true,
  linkcolor=RoyalBlue,
  urlcolor=RoyalBlue,
  citecolor=RoyalBlue,
  anchorcolor=RoyalBlue
}
\settowidth{\cftsecnumwidth}{Chapter 9: }
\titleformat{\section}
  {\normalfont\Large\bfseries}
  {Chapter \thesection:}
  {0.5em}
  {}
% \titleformat{\section}{
%     \scshape\raggedright\large\bfseries
%     }{}{0em}{}[\color{black}\titlerule]
%     \titlespacing{\section}{0pt}{5pt}{5pt}

\setstretch{1.5}
\setlist[itemize]{topsep=4pt, itemsep=3pt, parsep=0pt, partopsep=0pt, leftmargin=*}
\setlist[enumerate]{topsep=4pt, itemsep=3pt, parsep=0pt, partopsep=0pt, leftmargin=*}
\setlength{\footnotesep}{15pt} 
\setlength{\skip\footins}{15pt} 
\setlength{\parskip}{5pt}

\begin{document}

{\centering
  {\scshape \textls[50]{\fontsize{25}{30}\selectfont Reza {\bfseries Aliasgari Renani}}} \\ 
  {\fontsize{10}{10}\selectfont 08/01/2026} $|$
  {\fontsize{10}{10}\selectfont Motivation Letter} $|$
  {\fontsize{10}{10}\selectfont Dresden University of Technology} $|$
  {\fontsize{10}{10}\selectfont PhD (E-Sailors 6)} $|$
  {\fontsize{10}{10}\selectfont Electron Emitter}
\par}
\vspace{-20pt}
\noindent\rule{\textwidth}{0.5pt}

\vspace{-10pt}

My research objective is to contribute to E-sail technology by developing reliable high-voltage subsystems for nanospacecraft. The E-Sailor 6 doctoral position at TU Dresden focuses on a 20 kV electron emitter for ESTCube-LuNa under strict mass and volume limits while accounting for plasma and atomic-oxygen effects, with the electrical interface co-developed with E-Sailor 5. This area matches my background in plasma physics, electron-beam environments, and hardware testing. I am completing the final year of my Master's degree in Plasma Physics at Moscow Institute of Physics and Technology, where my work on radiation effects and FPGA systems has motivated me to pursue a PhD in space systems, focused on reliable propulsion hardware for missions like ESTCube-LuNa.

I joined Dr.\ Koveshnikov's laboratory in the Russian Academy of Sciences over two years ago to study irradiated semiconductor devices. We investigated the effects of electron and gamma irradiation and hydrogen plasma treatments on \ce{SiO2}-based MOS devices, distinguishing trap contributions by location (oxide, semiconductor, interface) and carrier type (electrons, holes) using electrical characterization techniques. This work resulted in a first-author publication\footnote{R. Aliasgari Renani, O.A. Soltanovich, M.A. Knyazev, S.V. Koveshnikov, Investigation of low energy electron irradiated \ce{SiO2} based MOS devices by C-V and TSC techniques, \textit{Russian Microelectronics}, 2023, DOI: \href{https://doi.org/10.1134/S1063739723600516}{10.1134/S1063739723600516}} and strengthened my understanding of device stability under bias stress and radiation. I also developed optimized MATLAB automation programs\footnote{Reza Aliasgari Renani, \textit{Automation Programs for Electrical Characterization Techniques}, available online: \href{https://github.com/rezaaliasgarirenani/IMT-Automation}{https://github.com/rezaaliasgarirenani/IMT-Automation}} for the experimental techniques such as TSC, C-V, I-V, I-T, and DLTS measurements, which enabled us to conduct the experiments and reduced run-times.

After my undergraduate degree, I joined the System-on-Chip Development Laboratory under the supervision of Dr.\ Vasiliev to develop FPGA-based systems. I implemented video-processing algorithms in Verilog using both manual coding and HDL generated from Simulink models, built verification testbenches, and used Python for quality control. For these pipelines, I wrote a fixed-point library and LUT-based function approximations to support efficient DSP. I then synthesized the design, checked resource utilization (LUTs and BRAM), and validated it on a Xilinx Zybo Z7 with live camera input and HDMI output. The codebase was maintained under Git-based source control.

Concurrently, at the Laboratory of Plasma Systems, I conducted research on radiation and plasma effects on FPGA devices. We exposed an FPGA board to 25-60 keV electron beams in low-pressure oxygen, generating plasma and X-rays from a tungsten target, while applying thermal cycling. The FPGA output signal was monitored and preliminary functionality tests were performed.

I have worked on engineering projects\footnote{Project repositories: \href{https://github.com/icarus-imperium/rocket-2025}{Model Rocket (Icarus Imperium)}, \href{https://github.com/rezaaliasgarirenani/Rover}{Model Lunar Rover}, \href{https://github.com/rezaaliasgarirenani/Aircraft-Detection-System}{Aircraft Detection System}, \href{https://github.com/rezaaliasgarirenani/Non-Conservative-Electric-Fields-and-Voltmeters}{Non-Conservative Electric Fields and Voltmeters}} both in teams and independently. In group projects, I built model rockets (40 and 60 N·s impulse) and completed multiple launches, and I collaborated on a model lunar rover through several prototypes. I developed an aircraft detection system using photoresistors and transistor circuits with rotational and translational tracking. In an individual electronics experiment on non-conservative electric fields, I designed the setup, gathered data, and documented how voltmeter readings become path-dependent under a time-varying magnetic flux. I also designed a CubeSat deployer and vibration fixture and validated prototypes through vibration testing.

For the E-Sailor 6 PhD position, I intend to focus on the electron emitter design, lifetime testing, and integration constraints for ESTCube-LuNa, including the plasma and atomic-oxygen environment. My background with electron-beam exposure and reliability testing prepares me to evaluate emission performance, contamination risks, and degradation mechanisms under realistic conditions. I will perform comparisons of emitter designs, carry out vacuum testing, and define the interface with E-Sailor 5. I am prepared to contribute to payload integration at TU Dresden and during the Aalto secondment. I will document results in technical reports and publications and present progress to project partners. This position aligns with my goal to build a career that bridges academia and industry in space exploration, spacecraft systems, and nanosatellite technologies. It would be an honor to study and work at TU Dresden in Germany, a country with longstanding traditions of scientific and cultural excellence.


\end{document}
